%% LyX 2.5.1 created this file.  For more info, see https://www.lyx.org/.
%% Do not edit unless you really know what you are doing.
\documentclass[brazil]{article}
\usepackage[T1]{fontenc}
\usepackage[utf8]{luainputenc}
\usepackage{babel}
\usepackage{mathtools}
\usepackage{amsmath}
\usepackage{amssymb}
\usepackage{geometry}
\geometry{verbose,tmargin=3cm,bmargin=3cm,lmargin=3cm,rmargin=2.5cm}
\usepackage[bookmarks=false,
 breaklinks=false,pdfborder={0 0 1},backref=false,colorlinks=false]
 {hyperref}

\makeatletter
%%%%%%%%%%%%%%%%%%%%%%%%%%%%%% User specified LaTeX commands.
\usepackage{babel}
















\makeatother

\begin{document}
\title{Postulados}
\author{Jhordan Silveira de Borba}
\maketitle
\begin{quote}
There is no general consensus as to what its fundamental principles are, how it should be taught, or what it really “means.” Every competent physicist can “do” quantum mechanics, but the stories we tell ourselves about what we are doing are as various as the tales of Scheherazade, and almost as implausible (\cite{griffiths2017introduction}). 
\end{quote}

\section{Analogy with Classical Mechanics}

\subsection{Newton's Second Law}

A law of physics can be understood as a rule that describes---usually by means of a mathematical formulation---how certain physical systems behave(\cite{Siddiqui2014}). Newton's second law, for example, can be expressed by the equation

\begin{equation}
F\left(x,t\right)=m\frac{d^{2}}{dt^{2}}x\left(t\right)
\end{equation}

where $x\left(t\right)$ represents the position of a particle of mass $m$ relative to the origin at the instant $t$. When we solve this differential equation under certain initial conditions, we obtain as a solution a function $x\left(t\right)$. Thus, once the forces acting on the particle and its initial conditions are specified, Newton's second law determines its temporal evolution: solving the equation provides us with the equation of motion of the particle.

For example, for a free particle($F=0$), we can obtain its equation of motion from Newton's second law:

\begin{equation}
0=m\frac{d^{2}x}{dt^{2}}
\end{equation}

Defining the velocity as $v\left(t\right)\equiv\frac{dx}{dt}$, we have:

\begin{align}
0 & =\frac{d^{2}x}{dt^{2}}\\
0 & =\frac{dv}{dt}\\
0\cdot\int^{t}_{t_{0}}dt' & =\int^{t}_{t_{0}}\frac{dv}{dt'}dt'\\
0 & =v\left(t\right)-v\left(t_{0}\right)
\end{align}

Therefore, $v\left(t\right)=v\left(t_{0}\right)$, that is, the particle's velocity is constant. Denoting this constant velocity by $v$, we have:

\begin{align}
v\left(t\right) & =v\left(t_{0}\right)\\
\frac{dx}{dt} & =v\\
\int^{t}_{t_{0}}\frac{dx}{dt'}dt' & =v\int^{t}_{t_{0}}dt'\\
x\left(t\right)-x\left(t_{0}\right) & =v\left(t-t_{0}\right)
\end{align}

If $x\left(t_{0}\right)=x_{0}$ and $t_{0}=0$, we obtain:

\begin{equation}
x\left(t\right)=x_{0}+vt
\end{equation}

In this case, $x\left(t\right)$ is easily interpretable as the position of the particle at the instant $t$.

\subsection{Schrödinger Equation}

The quantum mechanical analogue of Newton's second law is the Schrödinger equation. In one dimension and in the non-relativistic case, the Schrödinger equation is:

\begin{equation}
i\hbar\frac{\partial}{\partial t}\Psi\left(x,t\right)=-\frac{\hbar^{2}}{2m}\frac{\partial^{2}}{\partial x^{2}}\Psi\left(x,t\right)+V\left(x,t\right)\Psi\left(x,t\right)
\end{equation}

The corresponding equation for a free particle ($V\left(x,t\right)=0$) is:

\begin{align}
i\hbar\frac{\partial\Psi}{\partial t} & =-\frac{\hbar^{2}}{2m}\frac{\partial^{2}\Psi}{\partial x^{2}}
\end{align}
Using the method of separation of variables, that is, assuming that the solution can be separated into a spatial part $\Phi\left(x\right)$ and a temporal part $\Pi\left(t\right)$, that is, $\Psi\left(x,t\right)=\Phi\left(x\right)\Pi\left(t\right)$. Then:

\begin{align}
i\hbar\frac{\partial}{\partial t}\left(\Phi\left(x\right)\Pi\left(t\right)\right) & =-\frac{\hbar^{2}}{2m}\frac{\partial^{2}}{\partial x^{2}}\left(\Phi\left(x\right)\Pi\left(t\right)\right)\\
i\hbar\frac{1}{\Pi}\frac{\partial\Pi}{\partial t} & =-\frac{\hbar^{2}}{2m}\frac{1}{\Phi}\frac{\partial^{2}\Phi}{\partial x^{2}}
\end{align}

Since one side depends only on the spatial dimension ($x$) and the other only on the temporal dimension ($t$), the only way this can be true is if both sides are equal to a constant: $E$. Solving then for the temporal part:

\begin{align}
\frac{i\hbar}{\Pi}\frac{\partial\Pi}{\partial t} & =E\\
\frac{d\Pi}{dt} & =-\frac{iE}{\hbar}\Pi
\end{align}

Since the derivative of the function $\Pi\left(t\right)$ is equal to the function itself multiplied by a constant, we have an exponential solution $\Pi\left(t\right)=A\exp\left(-\frac{iE}{\hbar}t\right)$. For the spatial part:

\begin{align}
\frac{1}{\Phi}\frac{\partial^{2}\Phi}{\partial x^{2}} & =-\frac{2mE}{\hbar^{2}}\\
\frac{\partial^{2}\Phi}{\partial x^{2}} & =-\frac{2mE}{\hbar^{2}}\Phi
\end{align}

Since this is a second-order linear homogeneous differential equation with constant coefficients, we can look for solutions of the form $\Phi\left(x\right)=e^{\pm ikx}$. In this case, $e^{ikx}$ and $e^{-ikx}$ are two linearly independent solutions to $\frac{d^{2}\Phi}{dx^{2}}=-k^{2}\Phi$. By the linearity of the equation, its general solution is a linear combination of these solutions:

\begin{equation}
\Psi\left(x,t\right)=\Pi\left(t\right)\Phi\left(x\right)=A\exp\left(-\frac{iE}{\hbar}t\right)\left(B\exp\left(-i\sqrt{\frac{2mE}{\hbar^{2}}}x\right)+C\exp\left(i\sqrt{\frac{2mE}{\hbar^{2}}}x\right)\right)
\end{equation}

Or, writing $k^{2}=2mE/\hbar^{2}$ and $\omega=\frac{E}{\hbar}$:

\begin{align}
\Psi\left(x,t\right) & =K_{1}e^{-i\left(kx+\omega t\right)}+K_{2}e^{i\left(kx-\omega t\right)}
\end{align}
where $K_{1}$ and $K_{2}$ are constants determined by the conditions imposed on the solution. Here, some questions already arise. When we solved Newton's second law, we obtained $x\left(t\right)$, which described the position of the particle. When we solve the Schrödinger equation, we obtain $\Psi\left(x,t\right)$, but what is $\Psi\left(x,t\right)$? What does $\Psi\left(x,t\right)$ describe?

A somewhat more in-depth discussion of this can be found \href{/dupla}{here}.

\section{Postulates of quantum mechanics}

\subsection{Notation}

To move forward, we need to know the postulates\footnote{A postulate is a proposition assumed to be true within a theory, serving as a starting point for the construction of its concepts and results.} of quantum mechanics. The first point to note is that they can be formulated in different ways from one book to another. We will use Professor Jonas Maziero's material (\cite{maziero}). We need to begin, then, with some comments on Dirac notation.

In this notation, we write a vector as $\left|v\right\rangle $, called a ket. Taking the conjugate transpose of this vector (or adjoint), we obtain $\left\langle v\right|=\left|v\right\rangle ^{\dagger}=\left(\left|v\right\rangle ^{T}\right)^{*}$, which is called a bra. If we represent the vector $\left|v\right\rangle $ in an orthonormal basis $\left\{ \left|b_{j}\right\rangle \right\} $, we can write it as a linear combination of the basis vectors: $\left|v\right\rangle =\sum_{i}\left\langle b_{i}|v\right\rangle \left|b_{i}\right\rangle $. In matrix representation, we have

\begin{equation}
\left|v\right\rangle =\left[\begin{array}{c}
\left\langle b_{1}|v\right\rangle \\
\vdots\\
\left\langle b_{n}|v\right\rangle 
\end{array}\right]
\end{equation}
The product of a bra and a ket is called the inner product and results in a scalar:

\begin{equation}
\left\langle w|v\right\rangle =\left|w\right\rangle ^{\dagger}\left|v\right\rangle =\left[\begin{array}{ccc}
\left\langle b_{1}|w\right\rangle ^{*} & \dots & \left\langle b_{n}|w\right\rangle ^{*}\end{array}\right]\left[\begin{array}{c}
\left\langle b_{1}|v\right\rangle \\
\vdots\\
\left\langle b_{n}|v\right\rangle 
\end{array}\right]=\sum^{n}_{j=1}\left\langle b_{j}|w\right\rangle ^{*}\left\langle b_{j}|v\right\rangle 
\end{equation}
On the other hand, the product of a ket and a bra is called the outer product and results in an operator, whose matrix representation is:

\begin{equation}
\left|v\right\rangle \left\langle w\right|=\left|v\right\rangle \left|w\right\rangle ^{\dagger}=\left[\begin{array}{c}
\left\langle b_{1}|v\right\rangle \\
\vdots\\
\left\langle b_{n}|v\right\rangle 
\end{array}\right]\left[\begin{array}{ccc}
\left\langle b_{1}|w\right\rangle ^{*} & \dots & \left\langle b_{n}|w\right\rangle ^{*}\end{array}\right]=\left[\begin{array}{ccc}
\left\langle b_{1}|v\right\rangle \left\langle b_{1}|w\right\rangle ^{*} & \dots & \left\langle b_{1}|v\right\rangle \left\langle b_{n}|w\right\rangle ^{*}\\
\vdots & \vdots & \vdots\\
\left\langle b_{n}|v\right\rangle \left\langle b_{1}|w\right\rangle ^{*} & \dots & \left\langle b_{n}|v\right\rangle \left\langle b_{n}|w\right\rangle ^{*}
\end{array}\right]
\end{equation}

With this in mind, we can now move on to the postulates of quantum mechanics.

\section{State Postulate}
\begin{quote}
In every physical theory, we seek to describe the possible states of the system of interest, the changes in these states (temporal evolution), and the observable quantities and their measurements. For example, in Newtonian Mechanics (NM), the state of a system at a certain instant of time is specified by the values of position and velocity relative to a certain inertial reference frame, and its temporal evolution is given by Newton's second law. The observables are essentially the positions, from which we obtain the velocities using the time parameter. It is implicit in NM that the values of the observables are determined even before they are measured, and that all these observables can have well-defined values at a given instant of time. Next, we will see the postulates of QM for describing states, changes in these states, observables, and their measurements. We will see that the measurement part differs greatly from that of NM, mainly because of the existence of observables that cannot have well-defined values at the same time (\cite{maziero}). 
\end{quote}
The state of a quantum system is described by a normalized vector in a Hilbert space $\mathcal{H}$. For simplicity, we will consider $\mathcal{H}$ as a vector space equipped with an inner product function $\left\langle \cdot|\cdot\right\rangle :\mathcal{H}\rightarrow\mathbb{C}$.

Furthermore, a Hilbert space is a complete complex vector space. To understand what it means to be complete, we first need to understand the concept of a Cauchy sequence. A Cauchy sequence is a sequence whose elements become arbitrarily close to one another as we advance through the sequence. A complete metric space\footnote{A metric space is, essentially, a set of objects in which we define a way of measuring the distance between two objects. In the case of a vector space with an inner product, the distance can be defined from the norm induced by the inner product.} is one in which every Cauchy sequence converges to a limit that also belongs to the space itself.

By system, we refer to a degree of freedom\footnote{A degree of freedom corresponds to an independent physical variable that can vary and that can be used to characterize the state of a system. For example, the position and spin of a particle correspond to different degrees of freedom.}. For example, the spin of an electron is a system, another system would be its position, and so on. When we consider two or more degrees of freedom as a single entity, we call it a composite system, and if we divide the composite system into subsystems, its state is described by a normalized vector in a Hilbert space obtained by taking the tensor product of the individual spaces:

\begin{equation}
\mathcal{H}=\mathcal{H}_{a}\otimes\mathcal{H}_{b}\otimes\dots
\end{equation}


\section{Measurement Postulate}

Observable (measurable) quantities are represented by linear Hermitian operators, $O=O^{\dagger}$, defined on $\mathcal{H}$. When we measure an observable, we will obtain one of its possible values, which are given by the eigenvalues of the Hermitian operator.

Unlike classical mechanics (CM), quantum mechanics (QM), in general, does not allow us to predict with certainty the outcome of a measurement. What QM provides us with are the probabilities of obtaining the different outcomes, calculated by means of the Born rule:

\begin{equation}
P\left(o|\psi\right)=\left|\left\langle \psi|o\right\rangle \right|^{2}=\left|\left\langle o|\psi\right\rangle \right|^{2}
\end{equation}

This is the probability of obtaining the value $o$ when measuring the observable $O$ of a system that is in the state $\left|\psi\right\rangle $ immediately before the measurement, and $\left|o\right\rangle $ is the eigenvector of $O$ corresponding to the eigenvalue $o$. $O=O^{\dagger}$ are defined on $\mathcal{H}$. When we measure an observable, we obtain one of its possible values, which are represented by the eigenvalues of the Hermitian operator.

Before moving forward, we can stop to reflect for a moment. When we were working with classical mechanics and Newton's law, one of the observables of a particle was its position, and Newton's law directly described how this quantity varied in time.

However, when we approach quantum mechanics, the Schrödinger equation describes how the wave function varies in time (and space). With Dirac notation, as we will see when studying the dynamics postulate, we have a mathematical description of how the state vector of the system varies in time. However, the state vector is not itself an observable.

Even temporarily ignoring the probabilistic aspect, we use the state vector to calculate quantities associated with observables. This already constitutes a fundamental difference between the classical and quantum descriptions of systems. But if we calculate scalar quantities associated with observables from the state vector (or the wave function) of the system, what is the physical meaning of the state vector itself? Does it even have a direct physical meaning?

This is one of the points on which different interpretations of quantum mechanics diverge. There is still no consensus on which interpretation provides the correct description of physical reality, although some interpretations may be ruled out or considered problematic in light of certain theoretical arguments and experimental results.

Likewise, different interpretations diverge on the question of whether all relevant physical information about a system is necessarily contained in its state vector. We will explore different interpretations later on. For now, it is enough to emphasize that the physical meaning of the state vector and the status of the information it contains are profound interpretational questions on which there is no consensus.

\subsubsection*{Repeatability of measurements}

However, there is one more detail that, although it is not explicitly contained in the postulates we have presented so far, is significantly different when we are in the quantum domain: performing a measurement affects the system.

When we observe the position of a soccer ball that moves according to classical mechanics, we can, in principle, make this observation without significantly altering the state of the ball. We can therefore consider that the observation was performed approximately “externally” to the system itself, leaving it practically unchanged.

In the quantum domain, however, the situation is different. When we perform a measurement, the system interacts with the measurement apparatus and, in general, its state is significantly altered. One consequence of this is that, if we measure the same observable twice immediately in succession, the result of the second measurement will necessarily be equal to that obtained in the first.

Indeed, suppose that the first measurement of the observable ($O$) produced the result ($o$). After this measurement, the state of the system becomes a state associated with the result ($o$). In the case where the eigenvalue ($o$) is non-degenerate\footnote{A non-degenerate eigenvalue is one to which a single eigenvector corresponds, up to a global phase, that is, a complex phase factor $e^{i\phi}$ that multiplies the state vector and does not affect the measurement outcomes, since $e^{i\phi}\left(e^{i\phi}\right)^{*}=1$.}, we can write

\[
P\left(o|\psi\right)=1,\quad\rightarrow\left|\psi\right\rangle =e^{i\phi}\left|o\right\rangle 
\]

A second immediately subsequent measurement of the same observable will therefore once again produce the result ($o$).

Note that, theoretically, before the measurement is performed, quantum mechanics provides us only with the probabilities associated with the possible outcomes. In the laboratory, however, an individual realization of the measurement necessarily produces a definite outcome. This leads us to another fundamental question: the probabilistic nature of quantum mechanics. Its probabilistic predictions concern the frequencies and distributions of outcomes obtained from ensembles of systems prepared in equivalent ways and subjected to the same experimental procedure, and they do not, in general, determine which specific outcome will be obtained in an individual realization.

Before the measurement is performed, we have only the probabilities associated with the different values we may obtain. In this situation, in what state is the system? When dealing with classical mechanics, deterministic dynamics allows us, in principle, to determine the state of a system at an earlier instant $t'<t$ by knowing its state at an instant $t$. Can we extend this line of reasoning to quantum mechanics? Is the probabilistic character of the theory a fundamental feature of nature, or merely an expression of our ignorance about the state of the system? Again, we have no conclusive answers to this question, and different interpretations of quantum mechanics offer different answers.

And finally, still related to the probabilistic nature of measurements, misguided interpretations often support the idea that quantum mechanics is a scientific formulation that naturally incorporates the possibility of free will. However, the probabilistic character of the theory, by itself, provides no support for this conclusion and does not make quantum mechanics compatible with free will. We can discuss this misconception in greater detail when we address the dynamics postulate.

\subsubsection*{Quantum incompatibility}

Let $A$ and $B$ be two Hermitian operators defined on a finite-dimensional Hilbert space of dimension $d$. By the spectral theorem, we can write them in their respective spectral decompositions as

\begin{align}
A= & \sum^{d}_{j=1}a_{j}\left|a_{j}\right\rangle \left\langle a_{j}\right| & B= & \sum^{d}_{j=1}b_{j}\left|b_{j}\right\rangle \left\langle b_{j}\right|
\end{align}

The numbers $a_{j}$ and $b_{j}$ are the respective eigenvalues, and the sets$\left\{ \left|a_{j}\right\rangle \right\} ^{d}_{j=1}$ and $\left\{ \left|b_{j}\right\rangle \right\} ^{d}_{j=1}$ form orthonormal bases of eigenvectors of $A$ and $B$, respectively, that is, they diagonalize these operators, so that:

\begin{align}
A\left|a_{j}\right\rangle  & =a_{j}\left|a_{j}\right\rangle  & B\left|b_{j}\right\rangle = & b_{j}\left|b_{j}\right\rangle 
\end{align}

Assuming that the two operators share the same basis $\left|a_{j}\right\rangle =\left|b_{j}\right\rangle \coloneqq\left|j\right\rangle $, then, with the commutator defined as $\left[A,B\right]=AB-BA$, we have:

\begin{align}
\left[A,B\right] & =\left(\sum^{d}_{j=1}a_{j}\left|j\right\rangle \left\langle j\right|\right)\left(\sum^{d}_{k=1}b_{k}\left|k\right\rangle \left\langle k\right|\right)-\left(\sum^{d}_{k=1}b_{k}\left|k\right\rangle \left\langle k\right|\right)\left(\sum^{d}_{j=1}a_{j}\left|j\right\rangle \left\langle j\right|\right)\\
 & =\sum^{d}_{j,k=1}a_{j}b_{k}\left|j\right\rangle \left\langle j|k\right\rangle \left\langle k\right|-\sum^{d}_{j,k=1}b_{k}a_{j}\left|k\right\rangle \left\langle k|j\right\rangle \left\langle j\right|\\
 & =\sum^{d}_{j,k=1}a_{j}b_{k}\left|j\right\rangle \left\langle k\right|\delta_{jk}-\sum^{d}_{j,k=1}b_{k}a_{j}\left|k\right\rangle \left\langle j\right|\delta_{jk}\\
 & =\sum^{d}_{j=1}\left(a_{j}b_{j}\left|j\right\rangle \left\langle j\right|-b_{j}a_{j}\left|j\right\rangle \left\langle j\right|\right)\\
 & =0
\end{align}

Then:

\begin{equation}
0=\left[A,B\right]=AB-BA\Longrightarrow AB=BA
\end{equation}

Thus, we have that if the operators share the same set of eigenvectors, they commute. And to prove that if they commute, they share the same basis of eigenvectors:

\begin{align}
A\left(B\left|a_{j}\right\rangle \right) & =\left(AB\right)\left|a_{j}\right\rangle \\
 & =B\left(A\left|a_{j}\right\rangle \right)\\
 & =B\left(a_{j}\left|a_{j}\right\rangle \right)\\
 & =a_{j}B\left|a_{j}\right\rangle 
\end{align}

That is, in this case, we must have that $B\left|a_{j}\right\rangle $ and $\left|a_{j}\right\rangle $ belong to the eigenspace of $A$ corresponding to the eigenvalue $a_{j}$. If $a_{j}$ is non-degenerate, its eigenspace has dimension 1. Therefore, any eigenvector of $A$ associated with $a_{j}$ is linearly dependent on $\left|a_{j}\right\rangle $. Since $B\left|a_{j}\right\rangle $ also belongs to this eigenspace, we necessarily have $B\left|a_{j}\right\rangle =b_{j}\left|a_{j}\right\rangle $.

Now, if $a_{j}$ is degenerate, that is, if there are two or more linearly independent eigenvectors corresponding to the same eigenvalue $a_{j}$, then there exists a subspace of $\mathcal{H}$ generated by the eigenvectors of $A$ corresponding to $a_{j}$. It is a little more complex, but it is the same idea applied to a subspace of dimension greater than 1. To understand this, we will base our discussion on another source (\cite{MIT804_CommonEigenbases}).

For the degenerate case, we then have more vectors associated with $a_{j}$. Let us call this subspace $\mathcal{H}_{j}$ and, with $\left\{ \left|a^{k}_{j}\right\rangle \right\} ^{d_{j}}_{k=1}$ being any basis of $\mathcal{H}_{j}$, then $A\left|a^{k}_{j}\right\rangle =a_{j}\left|a^{k}_{j}\right\rangle $. Now, since $A$ and $B$ commute, we have $B\left|a_{j}\right\rangle \in\mathcal{H}_{j}$, that is, $B$ transforms vectors from $\mathcal{H}_{j}$ into vectors that also belong to $\mathcal{H}_{j}$. In particular, $B$ transforms an eigenvector of $A$ associated with $a_{j}$ into a vector that remains in the same eigenspace of $A$. If $B\left|a^{k}_{j}\right\rangle \neq0$, this vector is also an eigenvector of $A$ associated with $a_{j}$.

Therefore, each $B\left|a^{k}_{j}\right\rangle $ can be written as a linear combination of the basis vectors of $\mathcal{H}_{j}$. Thus, we can consider the action of $B$ “internally” within the subspace $\mathcal{H}_{j}$. Since $B$ is Hermitian, its restriction to $\mathcal{H}_{j}$ is also Hermitian. And since $\mathcal{H}_{j}$ has finite dimension (and is Hermitian), we can diagonalize this restriction of $B$ in an orthonormal basis of $\mathcal{H}_{j}$. Therefore, there exists a basis of $\mathcal{H}_{j}$, which we can denote by $\left\{ \left|a^{k}_{j}\right\rangle \right\} $, such that $B\left|a^{k}_{j}\right\rangle =b_{j}\left|a^{k}_{j}\right\rangle $. Since this new basis remains a basis of $\mathcal{H}_{j}$, its vectors continue to be eigenvectors of $A$ associated with $a_{j}$.

Thus, the eigenvectors are simultaneously eigenvectors of $A$ and $B$. Repeating this procedure for each eigenspace of $A$, we obtain an orthonormal basis of $\mathcal{H}$ consisting of common eigenvectors of the two operators.

In other words, if the Hermitian operators $A$ and $B$ commute, then they share the same orthonormal basis of eigenvectors. To better understand this proof, we can represent the operator $B$ as a matrix in the basis of eigenvectors of $A$. Using the completeness relation, $\mathbb{I}=\sum_{j}\left|a_{j}\right\rangle \left\langle a_{j}\right|$, we have $B=\mathbb{I}B\mathbb{I}=\sum_{j,k}\left(\left\langle a_{j}\left|B\right|a_{k}\right\rangle \right)\left(\left|a_{j}\right\rangle \left\langle a_{k}\right|\right)$.

Thus, the elements of the matrix representing $B$ in this basis are $B_{a_{j},a_{k}}=\left\langle a_{j}\left|B\right|a_{k}\right\rangle $, or in matrix form:

\begin{equation}
B=\left[\begin{array}{ccc}
\left\langle a_{1}\left|B\right|a_{1}\right\rangle  & \left\langle a_{1}\left|B\right|a_{2}\right\rangle  & \dots\\
\left\langle a_{2}\left|B\right|a_{1}\right\rangle  & \left\langle a_{2}\left|B\right|a_{2}\right\rangle  & \dots\\
\vdots & \vdots & \ddots
\end{array}\right]
\end{equation}

Then we can think of $B$ as a block-diagonal matrix:

\[
B=\left[\begin{array}{ccc}
B_{1} & 0 & 0\\
0 & B_{2} & 0\\
0 & 0 & \ddots
\end{array}\right]
\]

where each $B_{j}$ is a submatrix representing the internal action of $B$ acting on the subspace $\mathcal{H}_{j}$. If the eigenvalue $a_{j}$ is non-degenerate, then the block $B_{j}$ is just a $1\times1$ matrix. If $a_{j}$ is degenerate, then the block $B_{j}$ has dimension $d_{j}\times d_{j}$.

Since $B$ is Hermitian, each block $B_{j}$ is also a Hermitian matrix and, therefore, can be diagonalized. Thus, we can diagonalize $B$ within each subspace $\mathcal{H}_{j}$, obtaining a basis of eigenvectors of $B$ that are also eigenvectors of $A$. Thus, we prove that if and only if two operators $A$ and $B$ commute, then they share a common basis of eigenvectors. 
\begin{center}
\href{/exemplo-incompatibilidade}{Example 1}
\par\end{center}

In summary, when we have a non-degenerate eigenvalue $a_{j}$, the situation is very simple, since in this case, the only eigenvector of $A$, up to a multiplicative factor, is automatically also an eigenvector of $B$.

When $a_{j}$ is degenerate, its eigenvectors form a subspace $\mathcal{H}_{j}$, corresponding to the eigenvalue $a_{j}$, in which every nonzero vector is also an eigenvector of $A$ associated with $a_{j}$. Since $A$ and $B$ commute, applying $B$ to any vector in $\mathcal{H}_{j}$ results in another vector belonging to $\mathcal{H}_{j}$. Thus, it is sufficient to find, within this subspace, the vectors that are also eigenvectors of $B$.

This possibility is guaranteed by the fact that $B$ is Hermitian: its action restricted to $\mathcal{H}_{j}$ can be represented by a Hermitian submatrix, which can be diagonalized by an orthonormal basis of eigenvectors\footnote{The fact that a matrix is Hermitian guarantees that we can choose, in the space on which it acts, an orthonormal basis formed by its eigenvectors (\cite{DelftSpectralTheorem}, \href{https://interactivetextbooks.tudelft.nl/mqp-v2/operators.html}{on this page}.)}. Therefore, we can choose an orthonormal basis of $\mathcal{H}_{j}$ formed by eigenvectors of $B$. Since all vectors in $\mathcal{H}_{j}$ are also eigenvectors of $A$ associated with $a_{j}$, these vectors are common eigenvectors of $A$ and $B$.

Repeating this procedure for each eigenspace of $A$, we obtain an orthonormal basis of $\mathcal{H}$ formed by common eigenvectors of $A$ and $B$.

With the proof of this theorem, we know that if two observables $A$ and $B$ of a system do not commute, they cannot share all eigenvectors. If, by measuring the observable $A$, we prepare the system in a state that is not an eigenvector of the observable $B$, this vector must necessarily be a linear combination of the eigenvectors of with two or more nonzero probability amplitudes.
\begin{center}
\href{/exemplo-incompatibilidade}{Example 2}
\par\end{center}

\subsubsection*{Relação de Incerteza}

Aqui usaremos o formalismo cinemático quântico para obter relações matemáticas gerais que fornecem limitações para o produto ou soma das ``incertezas'' (desvios padrões) sobre dois observáveis.

\textbf{Relação de incerteza de Heisenberg-Robertson-Schrödinger (RIHRS)}

A RIHRS lê-se: 
\[
\text{Var}\left(A\right)\text{Var}\left(B\right)\ge(\text{CovQ}\left(A,B\right))^{2}+2^{-2}|\langle[A,B]\rangle|^{2},
\]
com a variância definida por $\mathrm{Var}(A)=\langle(A-\langle A\rangle)^{2}\rangle.$ As quantidades no lado direito da desigualdade serão definidas a seguir.

Para provar a RIHRS, usaremos como uma das principais ferramentas a desigualdade de Cauchy-Scharz (DCS): 
\begin{equation}
\langle\psi|\psi\rangle\langle\phi|\phi\rangle\ge|\langle\psi|\phi\rangle|^{2}=\langle\psi|\phi\rangle\langle\phi|\psi\rangle,
\end{equation}
com a igualdade obtida se e somente se $|\psi\rangle$ e $|\phi\rangle$ são colineares ($|\psi\rangle\propto|\phi\rangle$). Esse resultado é válido para qualquer par de vetores $|\psi\rangle,|\phi\rangle$ em um espaço de Hilbert $\mathcal{H}.$

Considere $A$ e $B$ dois observáveis de um sistema físico preparado no estado $|\xi\rangle$. Usaremos 
\[
\langle X\rangle:=\langle X\rangle_{\xi}=\langle\xi|X|\xi\rangle
\]
para denotar o valor médio de um operador qualquer $X$ e $I$ para denotar o operador identidade em $\mathcal{H}$. Antes de prosseguirmos, gostaria de exlorar mais esse resultado. Vamos considerar um sistema no estado $\left|\psi\right\rangle =\sum_{n}c_{n}\left|\psi\right\rangle $, e um operador $X$ com autovalores $x_{n}$ associados aos autovetores $\left|x_{n}\right\rangle $. Pela definição de média ponderada, o valor médio dos resultados é $\left\langle O\right\rangle =\sum_{n}o_{n}P\left(o_{n}\right)$, como $P\left(o_{n}\right)=\left|c_{n}\right|^{2}$, podemos escrever $\left\langle O\right\rangle =\sum_{n}o_{n}\left|c_{n}\right|^{2}$.

Como $\left\langle o_{n}|\psi\right\rangle =c_{n}$, então $\left\langle O\right\rangle =\sum_{n}o_{n}\left\langle \psi|o_{n}\right\rangle \left\langle o_{n}|\psi\right\rangle $ uma vez que $\left|\left\langle \psi|o_{n}\right\rangle \right|=\sqrt{\left\langle o_{n}|\psi\right\rangle \left\langle \psi|o_{n}\right\rangle }$. Então pela decomposição espectral $O=\sum_{n}o_{n}\left|o_{n}\right\rangle \left\langle o_{n}\right|$ então:

\[
\left\langle O\right\rangle =\left\langle \psi\right|\left(\sum_{n}o_{n}\left|o_{n}\right\rangle \left\langle o_{n}\right|\right)\left|\psi\right\rangle =\left\langle \psi\left|o\right|\psi\right\rangle 
\]

Vamos definir ainda os seguintes vetores: 
\begin{align}
|\psi\rangle & =(A-\langle A\rangle I)|\xi\rangle,\\
|\phi\rangle & =(B-\langle B\rangle I)|\xi\rangle,
\end{align}
que substituiremos na DCS. Notemos que 
\begin{align}
\langle\psi|\psi\rangle & =\left((A-\langle A\rangle I)|\xi\rangle,(A-\langle A\rangle I)|\xi\rangle\right)\\
 & =\left(|\xi\rangle,(A-\langle A\rangle I)^{\dagger}(A-\langle A\rangle I)|\xi\rangle\right)\\
 & =\langle\xi|(A-\langle A\rangle I)^{2}|\xi\rangle\\
 & =\mathrm{Var}(A),\\
\langle\phi|\phi\rangle & =\mathrm{Var}(B),
\end{align}
com $\mathrm{Var}(X)=\mathrm{Var}(X)_{\xi}=\langle\xi|(X-\langle X\rangle I)^{2}|\xi\rangle=\langle(X-\langle X\rangle I)^{2}\rangle$ sendo a variância de $X$. Também usamos $(A-\langle A\rangle I)^{\dagger}(A-\langle A\rangle I)=(A-\langle A\rangle I)^{2}$ uma vez que $A=A^{\dagger}$ e $\left\langle A\right\rangle =\left\langle A\right\rangle ^{*}$. Podemos ver que isso decorre de do fato de $A$ (e $I$) serem hermitianos uma vez que:

\begin{align*}
\left\langle A\right\rangle ^{*} & =\left(\left\langle \psi\left|A\right|\psi\right\rangle \right)^{*}=\left\langle \psi\left|A^{\dagger}\right|\psi\right\rangle =\left\langle \psi\left|A\right|\psi\right\rangle =\left\langle A\right\rangle 
\end{align*}

Ademais, temos que 
\begin{align}
\langle\psi|\phi\rangle & =((A-\langle A\rangle I)|\xi\rangle,(B-\langle B\rangle I)|\xi\rangle)\\
 & =(|\xi\rangle,(A-\langle A\rangle I)(B-\langle B\rangle I)|\xi\rangle)\\
 & =(|\xi\rangle,(AB-A\langle B\rangle-\langle A\rangle B+\langle A\rangle\langle B\rangle I)|\xi\rangle)\\
 & =\left\langle \xi\left|\left(AB-A\langle B\rangle-\langle A\rangle B+\langle A\rangle\langle B\rangle I\right)\right|\xi\right\rangle \\
 & =\left\langle \xi\left|AB\right|\xi\right\rangle -\left\langle \xi\left|A\right|\xi\right\rangle \left\langle \xi\left|B\right|\xi\right\rangle -\left\langle \xi\left|A\right|\xi\right\rangle \left\langle \xi\left|B\right|\xi\right\rangle +\left\langle \xi\left|A\right|\xi\right\rangle \left\langle \xi\left|BI\right|\xi\right\rangle \\
 & =\left\langle AB\right\rangle -\left\langle A\right\rangle \left\langle B\right\rangle -\left\langle A\right\rangle \left\langle B\right\rangle +\left\langle A\right\rangle \left\langle B\right\rangle \\
 & =\left\langle AB\right\rangle -\left\langle A\right\rangle \left\langle B\right\rangle 
\end{align}

Sendo: 
\begin{align}
\left[A,B\right] & =AB-BA,\\
\left\{ A,B\right\}  & =AB+BA
\end{align}
são, respectivamente, o comutador e anti-comutador de $A$ e $B$. Somando:

\begin{align*}
\left[A,B\right]+\left\{ A,B\right\}  & =AB-BA+AB+BA\\
 & =2AB\\
2^{-1}\left(\left[A,B\right]+\left\{ A,B\right\} \right) & =AB
\end{align*}
\[
AB=2^{-1}(\{A,B\}+[A,B]).
\]

Então:

\begin{align}
\left\langle \psi|\phi\right\rangle  & =\left\langle AB\right\rangle -\left\langle A\right\rangle \left\langle B\right\rangle \\
 & =\left\langle 2^{-1}(\{A,B\}+[A,B])\right\rangle -\left\langle A\right\rangle \left\langle B\right\rangle \\
 & =2^{-1}\left(\left\langle \{A,B\}\right\rangle +\left\langle [A,B]\right\rangle \right)-\left\langle A\right\rangle \left\langle B\right\rangle 
\end{align}

Até aqui temos então:

\begin{align}
\left\langle \psi|\psi\right\rangle \left\langle \phi|\phi\right\rangle  & \ge\left|\left\langle \psi|\phi\right\rangle \right|{}^{2}\\
\text{Var}\left(A\right)\text{Var}\left(B\right) & \geq\left|2^{-1}\left(\left\langle \{A,B\}\right\rangle +\left\langle [A,B]\right\rangle \right)-\left\langle A\right\rangle \left\langle B\right\rangle \right|{}^{2}
\end{align}

Como $\{A,B\}$ e $[A,B]$ são, respectivamente, operadores Hermitianos ($\{A,B\}^{\dagger}=\{A,B\}$) e anti-Hermitianos ($[A,B]^{\dagger}=-[A,B]$), seus valores médios são, respectivamente, números reais e imaginários puros. Que o valor médio de um hermitiano é um número real já verificamos, agora sobre um anti-hermitiano:

\begin{align*}
\left\langle A\right\rangle ^{*} & =\left(\left\langle \psi\left|A\right|\psi\right\rangle \right)^{*}=\left\langle \psi\left|A^{\dagger}\right|\psi\right\rangle =-\left\langle \psi\left|A\right|\psi\right\rangle =-\left\langle A\right\rangle 
\end{align*}

Agora se $\left\langle A\right\rangle =a+ib$, seu conjugado $\left\langle A\right\rangle ^{*}=a-ib$, se $\left\langle A\right\rangle ^{*}=-\left\langle A\right\rangle $, então $a+ib=-\left(a-ib\right)$, ou seja $a+ib=-a+ib$, pra isso ser igual então $a=0$. e ficamos apenas com $\left\langle A\right\rangle =ib$. Como $\left\langle A\right\rangle $ e $\left\langle B\right\rangle $ também são números reais. Podemos verificar sobre $\left[A,B\right]$ e $\left\{ A,B\right\} $

\begin{align*}
\left(A,B\right) & =AB\mp BA\\
\left(A,B\right)^{\dag} & =\left(AB\mp BA\right)^{\dagger}\\
 & =\left(AB\right)^{\dagger}\mp\left(BA\right)^{\dagger}\\
 & =B^{\dagger}A^{\dagger}\mp A^{\dagger}B^{\dagger}\\
 & =\mp A^{\dagger}B^{\dagger}+B^{\dagger}A^{\dagger}
\end{align*}

Ou seja, se $\left[A,B\right]=AB-BA$ então $\left[A,B\right]^{\dagger}=-\left[A,B\right]$. E se $\left\{ A,B\right\} =AB+BA$ então $\left\{ A,B\right\} ^{\dagger}=\left\{ A,B\right\} $. Então podemos escrever:

\begin{align*}
\left\langle \psi|\phi\right\rangle  & =\Re\left(\left\langle \psi|\phi\right\rangle \right)+i\Im\left(\left\langle \psi|\phi\right\rangle \right)\\
 & =\left[2^{-1}\left(\left\langle \left\{ A,B\right\} \right\rangle \right)-\left\langle A\right\rangle \left\langle B\right\rangle \right]+\left[2^{-1}\left\langle \left[A,B\right]\right\rangle \right]
\end{align*}

Ou seja:

\begin{align}
\Re(\langle\psi|\phi\rangle) & =2^{-1}\langle\{A,B\}-2\langle A\rangle\langle B\rangle\mathbb{I}\rangle,\\
\Im(\langle\psi|\phi\rangle) & =-i2^{-1}\langle[A,B]\rangle.
\end{align}

Obs.: Como $2^{-1}\langle[A,B]\rangle=ib$ é um nummero puramente imaginário, fazendo $-i\cdot\left(ib\right)=b$ obtemos seu coeficiente real. Assim, considerando que

\begin{align*}
\left|\left\langle \psi|\phi\right\rangle \right|^{2} & =\left\langle \psi|\phi\right\rangle ^{*}\left\langle \psi|\phi\right\rangle \\
 & =\left(\Re\left(\left\langle \psi|\phi\right\rangle \right)-i\Im\left(\left\langle \psi|\phi\right\rangle \right)\right)\left(\Re\left(\left\langle \psi|\phi\right\rangle \right)+i\Im\left(\left\langle \psi|\phi\right\rangle \right)\right)\\
 & =\Re\left(\left\langle \psi|\phi\right\rangle \right)^{2}+\Im\left(\left\langle \psi|\phi\right\rangle \right)^{2}
\end{align*}

Temos então:

\begin{align*}
\mathrm{Var}(A)\mathrm{Var}(B) & \geq(\Re(\langle\psi|\phi\rangle))^{2}+(\Im(\langle\psi|\phi\rangle))^{2}
\end{align*}

Agora vamos abrir: 
\[
\Im\left(\left\langle \psi|\phi\right\rangle \right)^{2}=\left(-i2^{-1}\langle[A,B]\rangle\right)^{2}=-2^{-2}\langle[A,B]\rangle^{2}
\]

Além disso como é puramente imginário então:

\[
\langle[A,B]\rangle^{2}=\left(-ib\right)^{2}=-b^{2}=-\left|b\right|^{2}=-\left|\left\langle \left[A,B\right]\right\rangle \right|^{2}
\]

Ficamos então com $\Im\left(\left\langle \psi|\phi\right\rangle \right)^{2}=2^{-2}\left|\left\langle \left[A,B\right]\right\rangle \right|^{2}$. E nossa desigualdade:

\begin{align*}
\left\langle \psi|\psi\right\rangle \left\langle \phi|\phi\right\rangle  & \ge\left|\left\langle \psi|\phi\right\rangle \right|{}^{2}\\
\text{Var}\left(A\right)\text{Var}\left(B\right) & \geq\Re\left(\left\langle \psi|\phi\right\rangle \right)^{2}+\Im\left(\left\langle \psi|\phi\right\rangle \right)^{2}\\
\text{Var}\left(A\right)\text{Var}\left(B\right) & \geq\Re\left(\left\langle \psi|\phi\right\rangle \right)^{2}+2^{-2}\left|\left\langle \left[A,B\right]\right\rangle \right|^{2}
\end{align*}

Nos resta agora demonstrar que:

\[
\Re\left(\left\langle \psi|\phi\right\rangle \right)^{2}=CovQ\left(A,B\right)^{2}=\left(\frac{Cov\left(A,B\right)+Cov\left(BA\right)}{2}\right)^{2}
\]

Onde a covariância é $\mathrm{Cov}(X,Y)=\langle XY\rangle-\langle X\rangle\langle Y\rangle$ entre os observáveis $X$ e $Y$, logo:

\begin{align*}
\Re\left(\left\langle \psi|\phi\right\rangle \right)^{2} & =\left(\frac{\langle AB\rangle-\langle A\rangle\langle B\rangle+\langle BA\rangle-\langle B\rangle\langle A\rangle}{2}\right)^{2}\\
\Re\left(\left\langle \psi|\phi\right\rangle \right)^{2} & =\left(\frac{\langle AB\rangle+\langle BA\rangle-2\langle B\rangle\langle A\rangle}{2}\right)^{2}\\
 & =\left(\frac{\left\langle AB+BA\right\rangle }{2}-\langle B\rangle\langle A\rangle\right)^{2}\\
 & =\left(2^{-1}\left\langle \left\{ AB\right\} \right\rangle -\langle B\rangle\langle A\rangle\right)^{2}
\end{align*}

E este $\Re\left(\left\langle \psi|\phi\right\rangle \right)=2^{-1}\left\langle \left\{ A,B\right\} \right\rangle -\left\langle A\right\rangle \left\langle B\right\rangle $ é o resultado que tínhamos obtido. Podemos escrever então $\Re\left(\left\langle \psi|\phi\right\rangle \right)^{2}=\text{CovQ}\left(A,B\right)^{2}$ e:

\begin{align*}
\text{Var}\left(A\right)\text{Var}\left(B\right) & \geq\Re\left(\left\langle \psi|\phi\right\rangle \right)^{2}+2^{-2}\left|\left\langle \left[A,B\right]\right\rangle \right|^{2}\\
\text{Var}\left(A\right)\text{Var}\left(B\right) & \geq\text{CovQ}\left(A,B\right)^{2}+2^{-2}\left|\left\langle \left[A,B\right]\right\rangle \right|^{2}
\end{align*}

Vale observar que se $A$ e $B$ são observáveis compatíveis, i.e., se $[A,B]=0$, teremos que $\mathrm{CovQ}(A,B)=\mathrm{Cov}(A,B)$.

Vamos analizar agora o problema da possível trivialidade da RIHRS, que tem por objetivo geral delimitar o quão pequeno pode ser o produto das variâncias de um par de observáveis. Para isso vamos assumir que o sistema está preparado em um estado que é autovetor do observável $A$, i.e., 
\[
|\xi\rangle=|a\rangle
\]
com $A|a\rangle=a|a\rangle$ e $a\in\mathbb{R}$. Neste caso pode-se verificar que 
\begin{equation}
\mathrm{Var}(A)_{|a\rangle}=\mathrm{CovQ}(A,B)_{|a\rangle}=\langle[A,B]\rangle_{|a\rangle}=0.
\end{equation}
E com isso a RIHRS nos retorna

\begin{equation}
0\mathrm{Var}(B)\ge0.
\end{equation}
Por conseguinte, neste caso, a RIHRS não fornece informação alguma sobre a possível incompatibilidade de $A$ e $B$. O mesmo tipo de limitação se aplica se $|\xi\rangle=|b\rangle$ para $B|b\rangle=b|b\rangle.$

De toda forma, vamos verificar. Para o comutador:

\begin{align*}
\left\langle a_{k}\left|\left[A,B\right]\right|a_{k}\right\rangle  & =\left\langle a_{k}\left|AB\right|a_{k}\right\rangle -\left\langle a_{k}\left|BA\right|a_{k}\right\rangle \\
 & =\left\langle a_{k}\left|\sum_{n}a_{n}\left|a_{n}\right\rangle \left\langle a_{n}\right|B\right|a_{k}\right\rangle -\left\langle a_{k}\left|B\sum_{n}a_{n}\left|a_{n}\right\rangle \left\langle a_{n}\right|\right|a_{k}\right\rangle \\
 & =\sum_{n}a_{n}\left[\left\langle a_{k}|a_{n}\right\rangle \left\langle a_{n}\left|B\right|a_{k}\right\rangle -\left\langle a_{k}\left|B\right|a_{n}\right\rangle \left\langle a_{n}|a_{k}\right\rangle \right]\\
 & =a_{k}\left[\left\langle a_{k}\left|B\right|a_{k}\right\rangle -\left\langle a_{k}\left|B\right|a_{k}\right\rangle \right]\\
 & =0
\end{align*}

Para a variância, vamos adotar outra estratégia, ao invés de escrever a decomposição espectral, vamos usar que $A|a\rangle=a|a\rangle$:

\begin{align*}
\mathrm{Var}(A)_{|a\rangle} & =\langle a_{n}|(A-\langle a_{n}\left|A\right|a_{n}\rangle I)^{2}|a_{n}\rangle\\
 & =\langle a_{n}|(A-a_{n}\langle a_{n}|a_{n}\rangle I)^{2}|a_{n}\rangle\\
 & =\langle a_{n}|(A-a_{n}I)^{2}|a_{n}\rangle\\
 & =\langle a_{n}|A^{2}|a_{n}\rangle+a^{2}_{n}\langle a_{n}|a_{n}\rangle-2a_{n}\langle a_{n}|A|a_{n}\rangle\\
 & =a^{2}_{n}+a^{2}_{n}-2a^{2}_{n}\\
 & =0
\end{align*}

E para a Covariância Quântica, já vimos que $\langle a_{n}\left|A\right|a_{n}\rangle=a_{n}$, $\left\langle a_{k}\left|AB\right|a_{k}\right\rangle =\left\langle a_{k}\left|BA\right|a_{k}\right\rangle =a_{k}\left\langle a_{k}\left|B\right|a_{k}\right\rangle $, então:

\begin{align*}
\mathrm{CovQ}(A,B)_{|a\rangle} & =\frac{\langle a_{k}\left|AB\right|a_{k}\rangle+\langle a_{k}\left|BA\right|a_{k}\rangle-2\langle a_{k}\left|B\right|a_{k}\rangle\langle a_{k}\left|A\right|a_{k}\rangle}{2}\\
 & =\frac{a_{k}\left\langle a_{k}\left|B\right|a_{k}\right\rangle +a_{k}\left\langle a_{k}\left|B\right|a_{k}\right\rangle -2a_{k}\langle a_{k}\left|B\right|a_{k}\rangle}{2}\\
 & =0
\end{align*}

\textbf{Relações de incerteza de Robertson e de Heisenberg (RIH)}

Ignorando-se\footnote{Isto é, se $a\geq b+c\geq c$, sendo $b\geq0$, então necessariamente $a\geq c$.} o termo relacionado à covariância quântica na RIHRS e tomando a raiz quadrada da expressão que resta, obtemos a relação de incerteza de Robertson:

\begin{align*}
\text{Var}\left(A\right)\text{Var}\left(B\right) & \geq\text{CovQ}\left(A,B\right)^{2}+2^{-2}\left|\left\langle \left[A,B\right]\right\rangle \right|^{2}\\
\sqrt{\text{Var}\left(A\right)}\sqrt{\text{Var}\left(B\right)} & \geq\sqrt{2^{-2}\left|\left\langle \left[A,B\right]\right\rangle \right|^{2}}\\
\sigma_{A}\sigma_{B} & \ge2^{-1}\left|\left\langle \left[A,B\right]\right\rangle \right|
\end{align*}

com e.g. $\sigma_{A}=\sqrt{Var(A)}.$ Quando aplicada a posição e momento linear é uma das relações de incerteza de Heisenberg: 
\begin{equation}
\sigma_{X}\sigma_{P}\ge2^{-1}|\langle[X,P]\rangle|=\hbar/2,
\end{equation}
pois $[X,P]=i\hbar I$.

Aqui estamos em condições de realizar mais algumas discussões. A relação de incerteza de Heisenberg é constantemente fruto de más compreensões. A primeira delas diz respeito ao próprio status da relação. Embora seja tradicionalmente chamada de princípio de incerteza, ela não precisa ser postulada como um princípio. Como vimos, a relação de Heisenberg pode ser obtida como um caso particular da Relação de incerteza de Heisenberg-Robertson-Schrödinger, que, por sua vez, pode ser derivada da desigualdade de Cauchy--Schwarz, juntamente com as propriedades matemáticas dos observáveis representados por operadores em um espaço de Hilbert.

O segundo ponto é que é incorreto interpretar a relação de incerteza como afirmando que uma partícula ``não tem posição e momento definidos'' ou que ``não podemos medir posição e momento na mesma preparação''. A primeira afirmação é uma interpretação possível, mas não é o conteúdo matemático da relação, enquanto a segunda é simplesmente incorreta. A relação de incerteza não diz que é impossível obter resultados de posição e momento na mesma preparação; ela estabelece uma relação entre as variâncias dos resultados dessas medidas.

Em outras palavras, podemos preparar repetidamente o sistema no mesmo estado e realizar medidas de posição e de momento. A partir dos resultados, podemos determinar as respectivas distribuições e, consequentemente, suas variâncias. A relação de incerteza estabelece que o produto dessas variâncias --- ou, mais precisamente, dos desvios-padrão --- deve satisfazer uma dada relação.

Isso contrasta com a mecânica clássica. Podemos imaginar, por exemplo, uma partícula parada em uma posição perfeitamente determinada. Ignorando limitações experimentais, seria possível conceber um estado clássico no qual a posição e a velocidade fossem simultaneamente bem definidas, por exemplo, uma particula parada na origem ($x=0,v=0$). Nesse caso, se repetíssemos as medidas, obteríamos sempre os mesmos valores e, portanto, teríamos $\sigma_{x}=\sigma_{v}=0$.

\subsection{Postulado da dinâmica}

As mudanças de estados em MQ são descritas por uma transformação unitária: 
\begin{equation}
|\psi_{t}\rangle=U_{t}|\psi_{0}\rangle,
\end{equation}
i.e., 
\[
U_{t}U^{\dagger}_{t}=U_{t}U^{\dagger}_{t}=\mathbb{I}.
\]
O operador de evolução temporal sendo obtido da equação de Schrödinger: 
\begin{equation}
i\hbar\partial_{t}U_{t}=H_{t}U_{t},
\end{equation}
com a condição inicial 
\[
U_{0}=\mathbb{I},
\]
onde $H_{t}$ é o operador Hamiltoniano do sistema no instante de tempo $t$,$\hbar$ é a constante de Planck dividida por $2\pi$ e usamos a notação $\partial_{t}=\frac{\partial}{\partial_{t}}$. As duas equações acima são equivalentes à equação de Schrödinger para o estado 
\begin{equation}
i\hbar\partial_{t}|\psi_{t}\rangle=H_{t}|\psi_{t}\rangle.
\end{equation}
Verificamos esse resultado como segue: 
\begin{align}
i\hbar\partial_{t}|\psi_{t}\rangle & =i\hbar\partial_{t}(U_{t}|\psi_{0}\rangle)\\
 & =(i\hbar\partial_{t}U_{t})|\psi_{0}\rangle\\
 & =H_{t}U_{t}|\psi_{0}\rangle\\
 & =H_{t}|\psi_{t}\rangle.
\end{align}

\textbf{Exemplo}: Quando $H_{t}=H$

Nesse tipo de situação, i.e., quando o Hamiltoniano não depende do tempo, podemos verificar que a seguinte forma geral é válida: 
\begin{equation}
U_{t}=e^{-iHt/\hbar}.
\end{equation}
Faremos isso primeiramente usando a decomposição espectral do Hamiltoniano, 
\[
H=\sum_{j}E_{j}|E_{j}\rangle\langle E_{j}|,
\]
e aplicando a função de operador: 
\begin{align}
i\hbar\partial_{t}U_{t} & =i\hbar\partial_{t}\exp\left(-\frac{iHt}{\hbar}\right)\\
 & =i\hbar\partial_{t}\exp\left(\sum_{j}\left(-\frac{iE_{j}t}{\hbar}\right)|E_{j}\rangle\langle E_{j}|\right)
\end{align}

Aqui devemos fazer uma pausa para entender a matemática. Vimos que para um operador hermitiano qualquer $T$ temos $\exp(T)=\sum^{\infty}_{n=0}\frac{1}{n!}T^{n}$. Primeiro podemos ver que $H^{2}$:

\begin{align*}
H^{2} & =\left(\sum_{j}E_{j}|E_{j}\rangle\langle E_{j}|\right)\left(\sum_{n}E_{n}|E_{n}\rangle\langle E_{n}|\right)=\sum_{j,n}E_{j}E_{n}|E_{j}\rangle\langle E_{j}|E_{n}\rangle\langle E_{n}|=\sum_{j}E^{2}_{j}|E_{j}\rangle\langle E_{j}|
\end{align*}

Onde usamos o fato de que $\left\langle E_{j}|E_{n}\right\rangle =\delta_{jn}$. De modo geral temos que $H^{n}=\sum_{j}E^{n}_{j}|E_{j}\rangle\langle E_{j}|$. Escrevendo então:

\begin{align*}
\exp(H) & =\sum^{\infty}_{n=0}\frac{1}{n!}H^{n}\\
 & =\sum^{\infty}_{n=0}\frac{1}{n!}\left(\sum_{j}E^{n}_{j}|E_{j}\rangle\langle E_{j}|\right)\\
 & =\sum_{j}\left(\sum^{\infty}_{n=0}\frac{E^{n}_{j}}{n!}\right)|E_{j}\rangle\langle E_{j}|
\end{align*}

E sendo $\exp(x)=\sum^{\infty}_{n=0}\frac{x^{n}}{n!}$, então: $\exp(H)=\sum_{j}e^{E_{j}}|E_{j}\rangle\langle E_{j}|$. No nosso caso temos então:

\begin{align}
i\hbar\partial_{t}U_{t} & =i\hbar\sum_{j}\partial_{t}\exp\left(-\frac{iE_{j}t}{\hbar}\right)|E_{j}\rangle\langle E_{j}|\\
 & =i\hbar\sum_{j}\left(-\frac{iE_{j}}{\hbar}\right)\exp\left(-\frac{iE_{j}t}{\hbar}\right)|E_{j}\rangle\langle E_{j}|\\
 & =i\hbar\left(-\frac{i}{\hbar}\right)\sum_{j}\exp\left(-\frac{iE_{j}t}{\hbar}\right)E_{j}|E_{j}\rangle\langle E_{j}|\\
 & =\sum_{j}e^{-\frac{iE_{j}t}{\hbar}}H|E_{j}\rangle\langle E_{j}\\
 & =H\left(\sum_{j}e^{-\frac{iE_{j}t}{\hbar}}|E_{j}\rangle\langle E_{j}|\right)
\end{align}
em que usamos $H|E_{j}\rangle=E_{j}|E_{j}\rangle$ ao longo de toda demonstração. as agora sendo então que 
\[
\sum_{j}e^{-\frac{iE_{j}t}{\hbar}}|E_{j}\rangle\langle E_{j}|=e^{-\frac{it}{\hbar}\sum_{j}E_{j}|E_{j}\rangle\langle E_{j}|}=e^{-\frac{iHt}{\hbar}|}
\]

Então $i\hbar\partial_{t}U_{t}=He^{-\frac{iHt}{\hbar}|}$ e como $i\hbar\partial_{t}U_{t}=H_{t}U_{t},$logo $U_{t}=e^{-iHt/\hbar}$.

OBS. Note que poderíamos obter $U_{t}$ resolvendo a equação de Schrödinger para cada hamiltoniano $H$. No entanto, tendo essa fórmula geral, podemos começar a resolução já calculando a exponencial, sem precisar resolver a equação diferencial.

Agora podemos abordar uma outra questão que se relaciona ao “determinismo” como oposição ao livre-arbítrio. Há pelo menos dois sentidos para a palavra determinismo. Um deles é mais associado à física clássica e consiste na ideia de que, se conhecemos perfeitamente o estado atual de um sistema e as leis que governam sua evolução, seu estado futuro e passado são, em princípio, determinados. O outro é um sentido mais amplo, de natureza filosófica, que diz respeito à negação do livre-arbítrio.

Quanto ao primeiro sentido, a formulação usual da mecânica quântica introduz um elemento probabilístico que contrasta com o determinismo da mecânica clássica. Os resultados das medidas não são, em geral, determinados de maneira unívoca pelo estado quântico anterior, mas descritos por distribuições de probabilidade. Nesse sentido, a teoria apresenta um elemento indeterminista. Vale notar, contudo, que a equação de Schrödinger que governa a evolução do estado quântico é determinista. Dado o estado inicial e o Hamiltoniano do sistema, a equação determina univocamente o estado quântico em qualquer instante posterior. As probabilidades dos resultados de uma medida são então obtidas a partir desse estado, por meio da regra de Born.

Se a natureza é, em última instância, determinista e a natureza probabilística da mecânica quântica decorre de alguma descrição incompleta, ou se a indeterminação quântica é fundamental, é uma questão para a qual não possuímos uma resposta consensual. Diferentes posicionamentos a respeito dessa questão dão origem, em parte, a diferentes interpretações da mecânica quântica.

Porém, é um equívoco pensar que esse elemento matemático probabilístico serviria, por si só, como suporte para uma defesa do livre-arbítrio. Não há, nas leis físicas conhecidas, uma formulação que incorpore o livre-arbítrio. Podemos aproximar a chance de obter cara ou coroa ao lançar uma moeda de maneira aproximadamente simétrica a 50\% para cada resultado. Esse é um evento probabilístico, mas não diríamos, por isso, que a moeda possui livre-arbítrio. A introdução de probabilidades na descrição da natureza, portanto, não confere mais livre-arbítrio à natureza do que o lançamento probabilístico de uma moeda.

Além disso, não apenas a evolução do estado quântico é determinista, como a própria distribuição de probabilidade em qualquer instante impõe restrições aos possíveis resultados e às suas probabilidades. Ao meu ver, essas restrições são incompatíveis com uma concepção de “livre” arbítrio na qual o agente deveria possuir liberdade irrestrita para determinar qualquer um dos resultados possíveis. 
\begin{quotation}
 \bibliographystyle{plain}
\bibliography{ref}
\end{quotation}

\end{document}
